\VerdierPhysicalPageAnchor{0031}%
\VerdierSemanticAnchor{intro-8-etale-lefschetz}
Dans le cadre de la topologie étale, on a une formule de Lefschetz analogue
[12].
\VerdierSemanticAnchor{intro-8-frobenius-L}
Cette formule, appliquée à l'endomorphisme de Frobenius d'une courbe
algébrique propre et lisse sur un corps fini, munie d'un faisceau
constructible pour la topologie étale, a permis la démonstration de la
rationalité des fonctions $L$ généralisées [11]. Signalons que dans cette
démonstration, on utilise une formule de Lefschetz où le faisceau des
coefficients possède effectivement des singularités.

\VerdierSemanticAnchor{intro-9}
\noindent\textbf{9.}\quad
Dans le cas de la formule de Lefschetz topologique, on a vu que le complexe :
\VerdierSemanticAnchor{intro-9-perfect-complex}
\[
\mathrm{RH}(X,F)
\]
est parfait, ce qui permet de définir le nombre de Lefschetz de :
\VerdierSemanticAnchor{intro-9-rh-phi}
\[
\mathrm{RH}(\Phi)\colon
\mathrm{RH}(X,F)\longrightarrow\mathrm{RH}(X,F).
\]

\VerdierSemanticAnchor{intro-9-cohomology-objects}
On voit là de façon particulièrement frappante pourquoi on ne peut remplacer
un complexe par la collection de ses objets de cohomologie. On ne peut en
effet définir la trace des endomorphismes :
\VerdierSemanticAnchor{intro-9-hi-phi}
\[
\mathrm{H}^i(\Phi)\colon
\mathrm{H}^i(X,F)\longrightarrow\mathrm{H}^i(X,F)
\]
que sous des conditions extrêmement restrictives.

\VerdierSemanticAnchor{intro-9-regular-trace}
Lorsque, par exemple, la dimension cohomologique globale de l'anneau $A$ est
finie, \emph{i.e.} lorsque $A$ est régulier, on peut définir la trace de
$\mathrm{H}^i(\Phi)$ en prenant des résolutions projectives finies. On a
alors :
\VerdierSemanticAnchor{intro-9-trace-formula}
\[
\chi\bigl(\mathrm{RH}(\Phi)\bigr)
=\sum_i(-1)^i\operatorname{Tr}\bigl(\mathrm{H}^i(\Phi)\bigr).
\]

\VerdierSemanticAnchor{intro-9-nonregular-rings}
Mais ainsi qu'on l'a vu, notamment en théorie des faisceaux pour la topologie
étale, on a à travailler avec des anneaux du type
$\mathbb{Z}/p^\ell\mathbb{Z}$ ($\ell\geq 2$) qui ne sont pas réguliers.

\VerdierSemanticAnchor{intro-9-k-theory}
Cette notion de complexe parfait est également utile en géométrie
algébrique : K-théorie et théorème de Riemann--Roch [14].

\VerdierSemanticAnchor{intro-10}
\noindent\textbf{10.}\quad
En conclusion, il semble qu'actuellement, le langage des catégories dérivées
soit l'outil indispensable pour permettre de formuler et de démontrer les
résultats essentiels des théories qu'on vient de mentionner.

\VerdierSemanticAnchor{intro-10-additive-scope}
Nous nous sommes limités dans ce travail à l'algèbre homologique additive :
catégories abéliennes et foncteurs additifs. Dès qu'on aborde les théorèmes
de Riemann--Roch [14] ce cadre est insuffisant.
\VerdierSemanticAnchor{intro-10-omissions}
Nous n'avons pas non plus abordé l'étude des différents groupes de classes,
ou groupes de Grothendieck, que l'on peut former avec les catégories
triangulées. Enfin, nous avons laissé de côté les questions d'homologie
relative, qui peuvent aussi se traiter dans le cadre des catégories
triangulées [15].

